% macros.tex --- Notation and shortcuts for the Collatz paper
% Authors: Alex Ye with Claude (Anthropic)

%--- Lagarias-convention operators ---
\DeclareMathOperator{\ord}{ord}
\DeclareMathOperator{\lcm}{lcm}
\DeclareMathOperator{\val}{v}             % 2-adic valuation
\DeclareMathOperator{\sgn}{sgn}

%--- The Collatz map and iterates ---
\newcommand{\T}{T}                          % Collatz / Syracuse map
\newcommand{\Tit}[1]{T^{#1}}                % T^k
\newcommand{\Cmap}{C}                       % accelerated Syracuse: C(n) = (3n+1)/2^{v_2(3n+1)}

%--- Stopping times ---
\newcommand{\stop}{\sigma}                  % stopping time (first n -> below n)
\newcommand{\totstop}{\sigma_\infty}        % total stopping time (first reach 1)
\newcommand{\stoptime}[1]{\stop(#1)}
\newcommand{\totstoptime}[1]{\totstop(#1)}

%--- Parity vector / Q-functions ---
\newcommand{\pv}{Q}                         % parity vector / Q-function
\newcommand{\pvk}[2]{Q_{#1}(#2)}            % first k entries of parity vector for n
\newcommand{\trajectory}[1]{\mathcal{O}(#1)}% forward orbit

%--- Density / measure ---
\newcommand{\logdens}{\overline{d_{\log}}}  % logarithmic density (upper)
\newcommand{\logdenslow}{\underline{d_{\log}}}
\newcommand{\nat}{\mathbb{N}}
\newcommand{\Z}{\mathbb{Z}}
\newcommand{\Zp}{\mathbb{Z}_p}              % p-adic integers
\newcommand{\Q}{\mathbb{Q}}
\newcommand{\R}{\mathbb{R}}

%--- Tao framework ---
\newcommand{\TaoConst}{c}                   % the implicit constant in Tao 2019
\newcommand{\logbound}{f(N)}                % function s.t. orbits go below f(N)

%--- Workflow markers ---
\newcommand{\TODO}[1]{\textcolor{red}{\textbf{[TODO: #1]}}}
\newcommand{\NOTE}[1]{\textcolor{blue}{\textbf{[NOTE: #1]}}}

%--- Misc shorthands ---
\newcommand{\eps}{\varepsilon}
\newcommand{\half}{\tfrac{1}{2}}
\newcommand{\bigO}{O}
\newcommand{\given}{\,|\,}
\newcommand{\E}{\mathbb{E}}
\newcommand{\Pr}{\mathbb{P}}
